\section{Discussion and Conclusion}
\label{sec:discussion}

\paragraph{Per-document choice.} A per-document chooser has to pick a
setup from a few questions per document, and the split-half oracle shows
what that choice is worth on the other questions: on the held-out
documents, nothing (\S\ref{sec:rq1}). The many ties explain why. When
most setups are within noise of the best, a choice made on half of a
document's questions is close to a random draw from the tied set, and a
larger choice set gives more room for unlucky draws. The decision holds
for each retriever and without the Indian documents, so it does not rest
on one retriever or on the questions whose validity we could not confirm.

\paragraph{The budget decides what ``best'' means.} The best chunk size
depends on how retrieval is scored (\S\ref{sec:rq3}). With a fixed
number of chunks, a larger chunk reads more text and so looks better; at
an equal token budget, smaller chunks win. Comparisons of chunkers at a
fixed $k$ therefore mix chunk size with the amount of text read, and we
recommend reporting results at an equal token budget.

\paragraph{Defaults and tuning.} The one large effect we found is a bad
default: LangChain's library default chunk size loses clearly to a size
tuned on other documents (\S\ref{sec:rq5}), while tuning our own
chunkers' sizes or the parser settings gives no significant gain.
Choosing a reasonable size once matters more than tuning per document.
For scanned PDFs the decisive step is detecting the scan, because a text
extractor reads nothing from it.

\paragraph{Answers.} Answer quality agrees with recall about the penalty
for large chunks but does not separate the small-chunk setups, and the
correlation between the two is weaker on the held-out documents
(\S\ref{sec:rq4}). Recall at a token budget is a measure of what the
reader gets to see; it does not guarantee a better answer.

\paragraph{Practical implication.} A document question-answering agent
can use one fixed setup with a small tuned chunk size and a scan check,
instead of selecting a chunker per document.

\paragraph{Conclusion.} In a pre-registered study with a frozen scorer
and held-out documents opened only after the plan was frozen, choosing a
chunking setup per document did not beat one fixed setup
(\RQoneDecision{}). Smaller chunks retrieve more at an equal token
budget, the usual fixed-$k$ measures reverse this ranking, and a poor
library default costs more than any per-document choice gains. We
publish the plan, the deviations and the result files so that the
analysis can be checked.
